\chapter{\introductionname}

En este primer capítulo presento la base conceptual y motivacional de mi Trabajo de Fin de Grado. Parto de una reflexión sobre el contraste de nuestra era tecnológica: un avance sin precedentes en comodidad y procesamiento de datos que no ha sido correspondido por una cultura de seguridad robusta. Planteo cómo, en un mundo donde cedemos nuestra privacidad de forma casi inercial, la Inteligencia Artificial surge no solo como un nuevo riesgo, sino como la herramienta definitiva para cerrar la brecha de inseguridad actual.

\section{La revolución de la Inteligencia Artificial}
En esta sección analizo la evolución reciente de la IA, explicando cómo hemos pasado de modelos de nicho a una democratización masiva liderada por herramientas como ChatGPT. Mi objetivo es exponer que, detrás de la interfaz de un chatbot, reside un motor de razonamiento ideal para trabajar con lenguajes técnicos. Argumentaré que el verdadero potencial de estos modelos no es responder preguntas, sino actuar como el núcleo de aplicaciones agénticas capaces de tomar decisiones de forma autónoma.

\section{El panorama de la ciberseguridad actual}
Aquí describo la crisis de seguridad informática que vivimos, caracterizada por la asimetría entre la velocidad del ataque y la lentitud de la defensa. Hablaré de la frecuencia alarmante de los \textit{data breaches} y de la exposición de datos sensibles. Me centraré en el problema de la falta de profesionales cualificados y en cómo la superficie de ataque (IoT, nube, sistemas distribuidos) crece mucho más rápido que nuestra capacidad humana para auditarla.

\section{Contexto y motivación personal}
En este apartado explico por qué este proyecto es el punto de convergencia de mis dos grandes pasiones: la IA y la ciberseguridad. Comento de forma honesta que, al no haber cursado asignaturas que profundizaran en el \textit{pentesting} avanzado o en arquitecturas agénticas al nivel que deseaba, utilizo este trabajo como mi propia vía de especialización académica. Presento este TFG como el puente necesario para mi futura carrera profesional, permitiéndome investigar de forma autónoma sobre la senda hacia la AGI aplicada a entornos reales.

\section{La unión: el Agente Autónomo de Seguridad}
Para finalizar la introducción, presento la visión idealista y técnica de mi propuesta. Describo el concepto de un auditor digital autónomo que unifica la potencia cognitiva de los modelos de lenguaje con el arsenal técnico de distribuciones como Kali Linux. Mi objetivo es plantear un sistema que no solo automatice escaneos, sino que sea capaz de trazar planes de auditoría complejos, interpretar resultados y aprender de los errores en tiempo real, convirtiéndose en un colaborador estratégico para el auditor humano.